\section{models@run.time Integration on GLSP}
\label{sec:models-run-time-glsp}

This section describes how \texttt{models@run.time} is integrated into the GLSP-based BP variant.
The integration is server-centric: the server owns the SSE connection, processes incoming events, and translates them into GLSP actions that are forwarded to the client.
The VS~Code extension contributes only the user-facing commands that start and stop the listener.
Once active, the listener filters received events, applies the necessary model updates, and dispatches corresponding actions to keep the diagram view synchronized with the running system.

\subsection{Server-Side Wiring}
\label{subsec:bp-server-wiring}

The BP server extends the generic UVL GLSP module and registers the SSE-related services as singletons.
The resulting control flow is straightforward: \texttt{FMBPEventListenerService} maintains the connection to the external FMBP event stream.
\texttt{FMBPHighlightActionDispatchService} and \texttt{FMBPContextUpdateService} independently consume parsed events and translate them into the appropriate GLSP action dispatches.

\lstinputlisting[
    language=java,
    float,
    floatplacement=ht,
    caption={ [\texttt{configureAdditionals}] Eager singleton bindings for SSE integration in the BP server.},
    label={lst:bp-server-sse-binding}
]{code/glsp_runtime_integration/bp_server_sse_binding.java}

\subsection{Initiating the Listener from VS~Code}
\label{subsec:sse-vscode-commands}

The SSE listener is controlled through three VS~Code commands contributed exclusively by the BP build profile: start listening, enable high-frequency polling, and stop listening.
Each command dispatches its corresponding GLSP action through the diagram connector, decoupling the UI trigger from the server-side implementation.

\lstinputlisting[
    float,
    floatplacement=ht,
    caption={ [\texttt{SSE commands}] SSE command contributions in the VS~Code extension.},
    label={lst:sse-command-contribution}
]{code/glsp_runtime_integration/sse_command_contribution.ts}

The commands are gated in the extension manifest so that they appear only when the BP-profile custom editor is active.
This prevents the SSE controls from being exposed to profiles that do not interact with FMBP, keeping the command palette uncluttered and the extension API surface minimal.

\subsection{SSE Listener and Event Parsing}
\label{subsec:sse-listener}

The network layer is implemented by \texttt{FMBPEventListenerService}.
On startup, the service waits for the FMBP health endpoint to become reachable before opening the event stream at \texttt{/events}.
The stream is consumed in a dedicated background thread, and the service automatically reconnects if the connection is lost.

\lstinputlisting[
    language=java,
    float,
    floatplacement=ht,
    caption={ [\texttt{FMBPEventListenerService}] Listener startup and stream consumption loop.},
    label={lst:fmbp-listener-loop}
]{code/glsp_runtime_integration/fmbp_listener_loop.java}

Each raw payload is parsed into a \texttt{ParsedServerSentEvent}, a normalized record that exposes the event type, source path, timestamp, and structured data map.
This representation abstracts over the raw JSON lines and allows downstream services to filter and route events without coupling them to the wire format.
Crucially, the \texttt{source} field identifies the feature model from which an event originates, allowing services to discard events that do not belong to the currently open diagram.

\lstinputlisting[
    language=java,
    float,
    floatplacement=ht,
    caption={ [\texttt{ParsedServerSentEvent}] Normalized SSE event record.},
    label={lst:parsed-sse-event}
]{code/glsp_runtime_integration/ParsedServerSentEvent.java}

\subsection{Translating Events into GLSP Actions}
\label{subsec:event-to-action}

The listener does not manipulate the diagram directly.
Instead, services register as data listeners on the \texttt{ServerSentEventsService} and are invoked selectively for the event types they handle.
This publisher-subscriber design keeps each service focused on a single concern and makes it straightforward to add support for new event types without modifying existing code.

\paragraph{Visual feedback via highlight actions.}
\texttt{FMBPHighlightActionDispatchService} subscribes to \texttt{requested}, \texttt{blocked}, and \texttt{waited\_for} events.
When such an event arrives, the service resolves the corresponding GLSP element identifiers from the b-thread and b-event name embedded in the payload, then dispatches a \texttt{HighlightElementAction} to the client.

\lstinputlisting[
    language=java,
    float,
    floatplacement=ht,
    caption={ [\texttt{FMBPHighlightActionDispatchService}] Highlight action dispatch for execution-state events.},
    label={lst:fmbp-highlight-dispatch}
]{code/glsp_runtime_integration/fmbp_highlight_dispatch.java}

This constitutes the primary visual feedback loop of the \texttt{models@run.time} integration.
The client receives these highlight actions and updates the diagram accordingly, providing an immediate visual indication of the active execution state without any polling.

\paragraph{Environment synchronization via context updates.}
\texttt{FMBPContextUpdateService} subscribes exclusively to \texttt{context\_update} events.
On receipt, the service compares the incoming attribute values against those held in the current model state, applies any changes, rebuilds the affected graphical node, and resubmits the updated model.
The client thereby receives a refreshed diagram representation reactively, again without polling.

\lstinputlisting[
    language=java,
    float,
    floatplacement=ht,
    caption={ [\texttt{FMBPContextUpdateService}] Incremental context update and model resubmission.},
    label={lst:fmbp-context-update}
]{code/glsp_runtime_integration/fmbp_context_update.java}

\subsection{Summary}
\label{subsec:models-runtime-summary}

The integration follows a well-defined pipeline.
A VS~Code command dispatches a GLSP action that instructs the server to open the SSE connection, representing the initial client-to-server leg of the flow.
From that point the data path reverses: the server parses each incoming FMBP event into a normalized record, which is forwarded to the registered data listeners.
Each listener maps the event to either a \texttt{HighlightElementAction}, which provides immediate visual feedback on the active execution state, or a model resubmission, which reflects updated environment attributes in the diagram.
The client receives both kinds of update through the existing GLSP action mechanism, requiring no additional communication channel and no polling on the client side.